\section{New deductions: the prescribed-measure loss as a concentration measure}
All assertions below concern the unchanged actual carrier and the entire
continuous class $\C$ described above. They are conditional only at the
explicit analytic-foundation boundary of that reconstruction. The prescribed
measure still ranges over every positive mass-two Radon $\mu$ with finite
full $B=D(\mu)$. Fix any of the fast actual maximizing sequences already
constructed. Its $a_n$, $O_m$, $N$, $u$, $e$, $p$, $r$, $J$ and $\nu$ retain
their original meanings. No operator is assigned to a new Borel function.

\begin{theorem}\label{fresh:concentration}
Modify $u$ on its $\mu$-null pole set only when forming integrals, and set
\[
 v_n=-a_n,\quad v=-u,\quad d_n=(v_n-v)_+,\quad k_n=(v-v_n)_+,
 \quad A=\bigcap_m\overline{O_m}.
\]
The positive Radon measures $d_n\mu$ have a nonempty weak cluster set.
Every cluster $\zeta$ has mass $J$, is supported on $A$, and satisfies
\[
 v_n\mu\ \longrightarrow\ v\mu+\zeta
\]
along the same convergent subnet. Conversely every weak cluster of
$v_n\mu$ has this form. In particular
\[
 J=0\quad\Longleftrightarrow\quad
 \|v_n\mu-v\mu\|_{\rm TV}\longrightarrow0
 \quad\Longleftrightarrow\quad v_n\mu\rightharpoonup v\mu.
\]
For every fixed $m$,
\[
 \lim_n\int_{O_m}d_n\,d\mu=J,
 \qquad \lim_n\int_{X\setminus O_m}d_n\,d\mu=0.
\]
There is no assertion that $\zeta$ is supported on $N$, is unique, is
singular to $\mu$, or avoids full-dimensional points in $A$.
\end{theorem}
\begin{proof}
The exact earlier $L^1$ calculation gives
\[
 \int k_n\,d\mu\to0,\quad \int d_n\,d\mu\to J,
 \quad \int v_n\,d\mu\to B-e=\int v\,d\mu+J.
\]
All densities here are Borel and integrable. Multiplication of a finite
Radon measure by an integrable Borel density gives a finite Radon measure:
truncate the density, approximate in $L^1$ by simple Borel densities,
and use regularity of the underlying finite measure. This proof applies
on compact Hausdorff $X$ and requires no metrizability.

We spell out the weak compactness being used. A sequence of positive
measures with bounded masses defines bounded coordinates
$\int h\,d\lambda_n$ for every $h\in C(X,\mathbb R)$. Tychonoff
compactness of the product of their closed bounded coordinate intervals
gives a convergent subnet. The limit coordinates preserve linearity and
positivity, and are bounded by the common mass bound times
$\|h\|_\infty$. Compact-Hausdorff Riesz representation gives a unique
positive Radon measure with those coordinates. This also proves
compactness of the entire weak cluster set. It is a net compactness
argument, not a measurable-net monotone convergence argument.
Apply it to $d_n\mu$; every cluster has mass $J$, by the constant test.
The identity
\[
 v_n\mu-v\mu=d_n\mu-k_n\mu,
 \qquad\|k_n\mu\|_{\rm TV}\to0
\]
proves both cluster-set statements.

For $n\ge m$, the uniform estimate on $X\setminus O_m$ gives
$d_n\le 2^{1-n}$ there. Since $\mu(X)=2$, its integral there is
at most $2^{2-n}$. Subtracting from the convergent total integral proves
both displayed localization limits. If a nonnegative continuous $h$ is
supported in $X\setminus\overline{O_m}$, its integral against
$d_n\mu$ tends to zero. The same holds for the cluster $\zeta$.
Urysohn's lemma and regularity imply
$\zeta(X\setminus\overline{O_m})=0$. The countable union of these
complements is $X\setminus A$, proving support on $A$.

The total variation norm in the theorem is exactly
$\int|a_n-u|\,d\mu\to J$. Thus $J=0$ implies both convergences.
Conversely weak convergence to $v\mu$, tested at $1$, implies
$B-e=\int v\,d\mu$, hence $J=0$. The support argument used closures
of open sets, and gives no stronger support assertion.
\end{proof}

\section{A sharp separation of bounded oscillation and deep negative tails}
\begin{theorem}\label{fresh:two-defects}
For $T>0$ put
\[
 b_n^T=(\min(v_n,T)-\min(v,T))_+,\quad
 q_n(T)=\int(v_n-T)_+\,d\mu,\quad
 q(T)=\int(v-T)_+\,d\mu.
\]
Then, for every fixed $T$,
\begin{equation}\label{fresh:sum}
 \int_N b_n^T\,d\mu+q_n(T)\longrightarrow J+q(T).
\end{equation}
Define
\[
 A_T^+=\limsup_n\int_N b_n^T\,d\mu,\quad
 A_T^-=\liminf_n\int_N b_n^T\,d\mu,\quad
 Q_T^+=\limsup_n q_n(T),\quad Q_T^-=\liminf_n q_n(T).
\]
Their limits as $T\to\infty$ exist and are finite; denote them
$A^\pm,Q^\pm$. They satisfy the exact identities
\[
 A^++Q^-=J,\qquad A^-+Q^+=J.
\]
If $\mu(N)=0$, then for every $T>0$ the ordinary scalar limit exists and
\begin{equation}\label{fresh:pure-tail}
 \lim_n q_n(T)=J+q(T),\qquad
 J=\lim_{T\to\infty}\lim_n\int(-a_n-T)_+\,d\mu.
\end{equation}
In this case $J=0$ is equivalent to uniform integrability of
$\{-a_n\}$ against $\mu$. Without the null-set antecedent, the exact
criterion is uniform integrability together with
$\min(v_n,T)\to\min(v,T)$ in $L^1(\mu)$ for every $T>0$.
Null-set annihilation alone is not the entire integral-tail criterion.
\end{theorem}
\begin{proof}
Let $c_n^T=(\min(v,T)-\min(v_n,T))_+$. The scalar truncation map is
nondecreasing and one-Lipschitz, so $0\le c_n^T\le k_n$ and its
integral tends to zero. Subtracting the truncated integrals gives
\[
 \int b_n^T\,d\mu-\int c_n^T\,d\mu
 =\int(v_n-v)\,d\mu-q_n(T)+q(T).
\]
The untruncated signed integral tends to $J$. On $X\setminus N$ we
have actual pointwise $v_n\to v$, hence $b_n^T\to0$. Since
$0\le b_n^T\le T$, dominated convergence on this measurable set
proves that its integral there tends to zero. This proves
\eqref{fresh:sum}.

For fixed nonnegative real $x,y$, the function
$(\min(x,T)-\min(y,T))_+$ is nondecreasing in $T$. Thus $A_T^\pm$
are nondecreasing. The tail quantities $Q_T^\pm$ are nonincreasing.
The preceding convergent sum, with nonnegative terms, bounds
$A_T^\pm\le J+q(T)$ and $Q_T^\pm\le J+q(T)$; the terms are finite
also for small $T$ because the integrals of $v_n$ are bounded.
For every $T$, elementary limsup/liminf of a sum converging to a real
number gives
\[
 A_T^++Q_T^-=J+q(T),\qquad A_T^-+Q_T^+=J+q(T).
\]
Sequential monotone convergence for the integrable $v$ gives
$q(T)\downarrow0$. Pass $T\to\infty$ to obtain the identities.
If $\mu(N)=0$, the bounded term in \eqref{fresh:sum} is zero, which
gives \eqref{fresh:pure-tail}.

For a sequence of integrable nonnegative functions, uniform
integrability is equivalent to
$\lim_{T\to\infty}\sup_n\int(v_n-T)_+\,d\mu=0$.
Indeed the usual tail integral on $\{v_n>2T\}$ is at most twice the
excess tail at $T$, and the excess tail is bounded by the usual tail.
Moreover this criterion is unchanged when $\sup_n$ is replaced by
$\limsup_n$: first control all sufficiently late $n$ at a fixed
threshold, and then increase the threshold to control the finitely
many remaining integrable functions. Hence uniform integrability is
equivalent to $Q^+=0$.
The truncated $L^1$ distance is
$\int(b_n^T+c_n^T)\,d\mu$. Its limit is zero for every $T$ exactly
when $A^+=0$, by the already proved convergence off $N$,
$\int c_n^T\to0$, and monotonicity in $T$. The two exact identities
show that $Q^+=0$ and $A^+=0$ together are equivalent to $J=0$.
They also show that $\mu(N)=0$ and uniform integrability imply $J=0$.
No assertion that full-dual finiteness supplies either antecedent has
entered this proof.
\end{proof}

\begin{proposition}\label{fresh:diagnostic}
Annihilation of an exceptional set and even full pointwise convergence
of continuous functions do not, as measure-theoretic assertions alone,
force the deep-tail defect to vanish. On $[0,1]$ with Lebesgue measure
there are nonnegative continuous $v_n$ with $v_n\to0$ at every point,
$\int v_n=1$, and $v_n\,dx\rightharpoonup\delta_0$.
Their exceptional open sets may have empty intersection while their
closures accumulate at $0$. This example is diagnostic only: its
functions are not asserted to be actual $\C$ weights or maximizing
competitors, and it is not a counterexample to P16.7.
\end{proposition}
\begin{proof}
For $n\ge2$, let $I_n=(2^{-n},2^{-n}+2^{-2n})$. These intervals are
pairwise disjoint. Let $v_n$ be the triangular continuous bump on
$\overline{I_n}$, zero at both endpoints, with central height
$2^{2n+1}$, and zero elsewhere. Its integral is one. Every point is
in the nonzero support of at most one bump, and $0$ is in none, so
pointwise convergence holds everywhere. For any continuous $h$,
\[
 \left|\int h v_n\,dx-h(0)\right|
 \le\sup_{x\in\overline{I_n}}|h(x)-h(0)|\to0.
\]
Set $O_m=\bigcup_{n\ge m}I_n$. Then $\bigcap_m O_m$ is empty,
but $\bigcap_m\overline{O_m}=\{0\}$. Thus the concentration
measure need not live on the intersection of the open sets. Finally
$\int(v_n-T)_+\to1$ for every fixed $T$, since
$\int\min(v_n,T)\le T2^{-2n}\to0$.
\end{proof}

\begin{proposition}\label{fresh:cluster-split}
Every concentration cluster $\zeta$ from
Theorem \ref{fresh:concentration} can be decomposed as
\[
 \zeta=\eta+\gamma,\quad \eta,\gamma\ge0,
 \quad \eta\ll\mu,\quad\eta(X\setminus N)=0,
 \quad\operatorname{supp}\gamma\subset A.
\]
For each integer $T\ge1$, along a further common subnet the bounded
parts $b_n^T\mu$ converge weakly to increasing measures $\eta_T$ with
$\eta_T\le T\mu|_N$ and $\eta_T\le\zeta$; $\eta=\lim_T\eta_T$.
The decomposition need not be unique, and no identification with
$A^\pm$ for the entire original sequence is asserted.
\end{proposition}
\begin{proof}
Apply the coordinate product compactness argument simultaneously to
$b_n^T\mu$ for the countably many integer $T$, refining a subnet on
which $d_n\mu\to\zeta$. The scalar inequalities $b_n^T\le b_n^{T+1}$
and $b_n^T\le d_n$ pass to nonnegative continuous tests and give the
order statements. At fixed $T$, the integral of $b_n^T$ off $N$ tends
to zero, while $b_n^T\le T$ everywhere. Consequently every
nonnegative continuous test has limiting integral at most its integral
against $T\mu|_N$. Order of finite Radon measures gives
$\eta_T\le T\mu|_N$.
The increasing limit $\eta(B)=\lim_T\eta_T(B)$ for Borel $B$ is a
finite measure, bounded by $\zeta$. It is Radon: approximate by
$\eta_T$ with total variation error
$\eta(X)-\eta_T(X)\to0$. It annihilates every $\mu$-null Borel set
and $X\setminus N$. The positive measure $\gamma=\zeta-\eta$ is
Radon and is supported on $A$, as is $\zeta$. Only countable
monotone measure operations were used.
\end{proof}

\section{The two scalar defects and the ordinary USC relaxation}
\begin{proposition}\label{fresh:JK}
For every eligible $\mu$ and every chosen fast sequence,
\[
 J-K=\int u\,d\mu-\int u\,d\nu=\int u\,d(\mu-\nu).
\]
Here both integrals are finite. Thus $J=K$ is exactly the equality of
these two Borel moments, not a consequence of the existing estimates.
The ordinary USC relaxation $R$ satisfies $u\le R\le0$, is
integrable against both $\mu$ and $\nu$, and obeys
\[
 e-\int R\,d\mu=B-J-\int(R-u)\,d\mu.
\]
The left side is only a scalar expression with the canonical number
$e$; it is not an assertion $E(R)=e$ or $R\in\mathcal G$.
\end{proposition}
\begin{proof}
The reconstructed tail estimates give $\int u\,d\nu=p$. The
definitions give $J=\int u\,d\mu-e+B$ and $K=p-e+B$; subtract.
Topological upper relaxation dominates the pointwise limsup, so
$u\le R\le0$. Integrability of $-u$ against the two measures
dominates $-R$. Substitution gives the last identity.
\end{proof}

\section{Bounded attainment removes the prescribed-measure loss}
\begin{theorem}\label{fresh:attainment}
Suppose an actual normalized $q\in\C$ attains the full dual of an
eligible $\mu$. Then, for every selected fast maximizing sequence,
\[
 I(a_n,q)\le4\varepsilon_n,\quad e=E(q),\quad
 \nu=\MA(q),\quad R=q,\quad J=0,
 \quad u=q\ \mu\hbox{-almost everywhere}.
\]
The same Borel equality holds almost everywhere for every bounded MA
and bounded mixed measure. The pointwise equality holds at every full
point. No bounded attainment for arbitrary $\mu$ is proved. In general
bounded attainment alone here does not identify $\mu$ with $\MA(q)$.
\end{theorem}
\begin{proof}
Apply the exact midpoint inequality to $a_n,q$, with the latter's
deficit zero. The energy and genuine weak-test moduli identify $e$
and $\nu$ as displayed. Actual atoms and the fixed-test absolute
estimates give $\ell(x)=q(x)$ at every full point. Substitute these
values into the all-affinoid Shilov formula. The classical maximum
principle gives $\rho(x)+q(x)\le\max_{\Gamma_A}(\rho+q)$ for every
affinoid neighborhood $A$; continuity and an affinoid neighborhood
inside each ordinary neighborhood give the converse after infimizing.
Thus $R=q$ on all $X$.
Since $u\le R=q$ and $u$ is $\mu$-integrable,
$\int u\,d\mu\le\int q\,d\mu=E(q)-B=r$. The earlier opposite
inequality $\int u\,d\mu\ge r$ forces equality, $J=0$, and
$u=q$ almost everywhere for $\mu$. For every fixed bounded MA test,
the $I$ modulus gives $a_n\to q$ in its $L^1$ space, while the fast
tail theorem gives $a_n\to u$ in the same space. Uniqueness of an
$L^1$ limit gives the claimed equality. Domination of a mixed measure
by twice a bounded midpoint MA measure gives its equality as well.
The argument proves no global equality at other points.
\end{proof}

\section{An equation-identifying slice using only actual invariant variations}
\begin{theorem}\label{fresh:toric-attainment}
Let $q(x)=\bar q(\tau(x))$ be a smooth periodic real weight such that
$I_2+\operatorname{Hess}\bar q$ is positive definite everywhere on
$\mathbb R^2/\mathbb Z^2$. If $q$ attains $D(\mu)$ for an eligible
arbitrary-support $\mu$, then $\mu=\MA(q)$. In particular, on the
entire eligible class,
\[
 D(\mu)=0\quad\Longleftrightarrow\quad\mu=\theta^2.
\]
The forward assertion forces skeletal support as a conclusion, not
as a restriction on the prescribed measure. This slice does not solve
the arbitrary-measure equation or analytic representation problem.
\end{theorem}
\begin{proof}
Set $\lambda=\MA(q)$. The coordinate tropical dimension rule and
identity calibration used for the reference apply to every smooth
coordinate weight: off the Gauss section the top form vanishes, and
on that section the raw measure is
$2\det(I_2+\operatorname{Hess}\bar q)\,dr_1dr_2$.
Thus $\lambda$ is supported on the Gauss section. The form is globally
defined under integer deck translations; no new geometric descent
assertion is required.

For every smooth periodic $h$, compactness and the positive definite
matrix give a uniform $c>0$ and a finite Hessian bound. Consequently
$q+th\circ\tau\in\C$ for both signs of all sufficiently small $t$:
the local potential is an ordinary smooth convex coordinate function.
First variation at the attained maximum gives
$\int h\circ\tau\,d(\lambda-\mu)=0$.
Smooth periodic tests are uniformly dense among continuous torus
tests, by periodic convolution with a positive smooth approximate
identity; uniform continuity proves uniform approximation.
Thus $\tau_*\lambda=\tau_*\mu$ as actual Radon measures.

For any $a\in\C$, the reconstructed Gauss inequality gives
$a(x)\le a(\sigma(\tau(x)))$. The section is continuous, so its
restriction is a continuous torus function. Projection equality and
skeletal support therefore give $\int a\,d\mu\le\int a\,d\lambda$.
On the other hand, the actual segment $(1-t)q+ta$ is in $\C$ for
$0\le t\le1$. Its right first derivative at the maximum is at most
zero, giving
$\int(a-q)\,d(\lambda-\mu)\le0$. Since $q$ is a coordinate weight,
projection equality gives $\int q\,d(\lambda-\mu)=0$.
Combining the two inequalities yields equality of integrals for every
$a\in\C$. Density of its real span and equality of finite masses
give $\mu=\lambda$.

If $D(\mu)=0$, then $q=0$ attains the dual because $F_\mu(0)=0$;
its coordinate matrix is $I_2$. The preceding theorem gives
$\mu=\theta^2$. Conversely tangent concavity at zero gives
$D(\theta^2)=0$. This proves the equivalence, without using a
full-support measure or an arbitrary smooth positivity margin.
\end{proof}

\section{New residual status and remaining audit boundary}
The positive concentration measures, their closure localization,
the exact bounded-oscillation/deep-tail split, the finite Borel identity
between $J$ and $K$, the bounded-attainment loss theorem, and the
invariant-variation slice are new complete deductions at the explicitly
qualified bounded-calculus scope. Every prescribed measure remains
eligible on the original full class; no domination assumption is used
in these universal deductions. The conclusions with an antecedent
state it openly and do not supply it for arbitrary measures.

In particular a proof that every eligible measure annihilates $N$
would still need the uniform negative-tail integral estimate. Uniform
integrability alone would still need bounded-truncation convergence on
$N$. Neither has been proved from full-dual finiteness. Concentration
on the intersection of the closures of $O_m$ cannot be excluded by
the fact that no full point lies in $N$.

We do not prove global USC or original $\C/\mathcal G$ membership
of $u$, positivity of tail envelopes, $J=0$ or $K=0$ for every
measure, or the genuine prescribed-measure existence-and-uniqueness
theory. Nonalgebraization, coherent descent, deeper smooth partition,
finite-flat/intersection, graded reduction and curve comparison source
trees remain incomplete independent audit boundaries. Whole historical
retention and compilation data supply no correctness grant. The
unchanged original P16.7 remains OPEN; candidate is false.
